\chapter*{ABSTRACT}
\begin{center}
\textbf{\ThesisTitle}\par
\ThesisAuthor
\end{center}

Dynamically balanced mobile robots combine a compact support polygon with active stabilization, but conventional two-wheel designs do not provide direct lateral motion. This thesis documents the design and development of Zephyr, a self-balancing mobile robot whose four Mecanum wheels are arranged along one common wheel-axis line rather than at the corners of a rectangular chassis. The project integrates four CubeMars AK40-10 actuators, an ISM330DHCX inertial measurement unit, a controller-area-network actuator bus, radio control, a local display, a six-cell lithium-polymer power source, embedded real-time firmware, and a developing custom printed-circuit board. A geometry-parameterized Mecanum-wheel Jacobian is derived so the model remains valid until the physical wheel order, roller handedness, and axial locations are measured. The balancing dynamics are represented by a nonlinear wheeled inverted-pendulum model and a linearized upright model suitable for state-feedback design. The current firmware demonstrates sensor acquisition, attitude-estimation infrastructure, radio-command decoding, display operation, actuator communication, feedback monitoring, and safety-gated command flow on an ESP32 prototype; the final balance controller, wheel-torque allocation, STM32 migration, and complete experimental validation remain in progress. Validation is therefore specified through reproducible simulation and hardware procedures rather than claimed as completed. \textless INSERT RESULT: Report final RMS pitch error, maximum pitch excursion, settling time, overshoot, steady-state error, disturbance-recovery limit, planar-velocity tracking error, control-loop timing, and CAN error statistics after testing.\textgreater{} The principal contribution of the present work is a traceable mechatronic architecture and analysis framework for evaluating whether a collinear four-wheel Mecanum arrangement can combine upright stabilization with controllable longitudinal, lateral, and yaw motion. Quantitative performance conclusions remain pending the measurements identified in this draft.

\noindent\textbf{Keywords:} self-balancing robot, Mecanum wheel, omnidirectional mobility, wheeled inverted pendulum, embedded control, CAN bus, mechatronics
\clearpage
